\chapter{总结与展望}

本项目在 xv6-RISC-V 原有调度框架内完成了四部分改造。进程控制块增加了优先级和等待计数字段，fork 可以安全继承父进程优先级。调度器按静态优先级选择可运行进程，并用轮转起点维持同级进程的基本公平。Aging 负责逐级提升持续等待的低优先级进程，系统调用和 \texttt{nice} 命令提供用户态查询与修改入口。实现沿用 xv6 的进程锁和状态转换，没有加入课题范围外的复杂调度结构。

选出优先级最高的进程只是调度修改的一部分。字段何时初始化和清零、fork 读取字段时是否持锁、多 CPU 下候选状态是否变化、Aging 如何衔接手动设置，都会影响内核运行。项目把这些约定落到少量字段、命名常量和明确的重置时机中，模块接口因而较容易核对。

第五阶段的专项测试程序已经合并到主分支，并在单 CPU、3 CPU 的 QEMU 环境中完成回归。\texttt{prioritytest} 覆盖了默认值、边界与非法参数、PID 查询与设置、\texttt{nice}、fork 继承、同级公平和不同优先级；\texttt{agingtest} 验证了高优先级竞争存在时低优先级进程能够逐级提升；两组配置下的 \texttt{usertests -q} 也均输出了 \texttt{ALL TESTS PASSED}。测试结果保存为原始日志与结构化结果文件，报告中的数值均来自本次实际 QEMU 执行。

本次集成还修复了测试运行器在 macOS 上读取 Linux 环境信息的问题，并在 \texttt{docs/bugs.md} 中记录了原因和验证方式。当前实现仍未把调度次数和运行时间作为新的持久化字段加入进程控制块；现阶段的调度统计由专项测试输出的 \texttt{METRIC}、\texttt{AGING\_RESULT} 以及运行器汇总文件承担。后续若课程需要更细粒度的性能分析，可在不改变现有调度语义的前提下增加专门统计字段和长期实验设计。
